\documentclass[a4paper]{article}
\usepackage[pdftex]{graphicx}
\usepackage[utf8]{inputenc}
\usepackage{enumerate}
\usepackage{amssymb}
\usepackage{icomma}
\usepackage{siunitx}
\sisetup{locale=DE}
\usepackage{tikz}
\usepackage{href-ul}
\hypersetup{
	colorlinks=true,
	linkcolor=blue,
	urlcolor=blue} 
\usepackage{geometry}
\geometry{a4paper, top=15mm, left=15mm, right=15mm, bottom=15mm,
	headsep=10mm, footskip=12mm}

\begin{document}
\begin{tikzpicture}(10,3)
	\draw[ultra thick](2,0) --(10,0) -- (10,1.3) --(2,1.3) -- (2,0);
	\draw[fill=black](2,0)-- (0,.6) -- (2,.6) -- (2,0);
	\node at (6,.6) {\large Lösungsblatt von \href{https://www.okuyakl.de}{www.okuyakl.de}};
\end{tikzpicture}
\vspace{0.5 cm}

\noindent{\bf Aufgabe 1. a)}\\
$$\mathbb{D}_f=\mathbb{R}^+ \qquad \lim\limits_{x\to\infty}x\cdot\ln x=\infty \qquad \lim\limits_{x\to 0}x\cdot\ln x=0 \quad \textnormal{vgl. Merkhilfe}$$
\vspace{0.5 cm}

\noindent{\bf Aufgabe 1. b)}\\
$$f(-x)=-x\cdot\ln (-x)\ne f(x)\ne-f(x)\quad \Rightarrow \quad \textnormal{keine Symmetrie}$$
\vspace{0.5 cm}

\noindent{\bf Aufgabe 1. c)}\\
$$
\renewcommand{\arraystretch}{2}
\begin{array}{rcll}

x\cdot\ln x&=&0  \\
x&=&0  &\not \in \mathbb{D}_f\\
\ln x&=&0 &\Rightarrow \quad x=1   \quad \textnormal{einfache Nullstelle}\\
\end{array}
$$
\vspace{0.5 cm}

\noindent{\bf Aufgabe 1. d)}\\
$$
\renewcommand{\arraystretch}{2}
\begin{array}{rcl}
f(x)&=&x\cdot\ln x  \\
f'(x)&=&\ln x +1 \\
f''(x)&=& {1\over x} \\
f'''(x)&=&-{1\over x^2}
\end{array}
$$
\vspace{0.5 cm}

\noindent{\bf Aufgabe 1. e)}\\
Wir bestimmen die Nullstellen der Ableitung
$$
\renewcommand{\arraystretch}{2}
\begin{array}{rcll}
f'(x)&=&\ln x +1&\stackrel{!}{=}0\\
-1&=&\ln x  &\Rightarrow \quad x=e^{-1} \\
\end{array}
$$
Monotonietabelle:
$$
\renewcommand{\arraystretch}{2}
\begin{array}{c|c|c|c|c|c}

& x<e^{-1} & x=e^{-1} &  x>e^{-1} \\
\hline
\ln x +1 & - & 0 & + \\
\hline
\hline
f'(x) & - & 0 & + \\
\hline
G_f & \searrow & \rightarrow & \nearrow \\
\hline
  & smf & T & sms & H & smf
\end{array}
$$
$$f(e^{-1})=-e^{-1} \quad \Rightarrow \quad \textnormal{Tiefpunkt} \quad T(e^{-1}|-e^{-1})$$
$G_f$ ist streng monoton fallend für $x \in ]0;e^{-1}]$\\
$G_f$ ist streng monoton steigend für $x \in [e^{-1};\infty]$
\newpage

\noindent{\bf Aufgabe 1. f)}\\
Es gibt keine Nullstellen der zweiten Ableitung wegen:
$${1\over x}> 0 \quad \textnormal{für alle}\quad  x \in \mathbb{D}_f$$

$G_f$ ist überall linksgekrümmt \\
\vspace{0.5 cm}

\noindent{\bf Aufgabe 1. g)}\\
Die Tangente hat die Form $y=mx+t$ mit:
$$x=1 \qquad y=0 \qquad m=f'(1)=1 $$
$$
\renewcommand{\arraystretch}{2}
\begin{array}{rcll}

y&=&mx+t \\
0&=&1\cdot 1 +t \\
t&=& -1 &\quad \Rightarrow \qquad y= x -1 

\end{array}
$$ 
\vspace{0.5 cm}

\noindent{\bf Aufgabe 1. h)}\\
$$F'(x)=x\cdot\left(\ln x -{1\over 2}\right) + {1\over2}x^2\cdot {1\over x}=x\cdot\ln x=f(x)$$

\begin{minipage}{0.5\textwidth}
		\noindent{\bf Aufgabe 1. i)}\\
		$$\int\limits_{u}^1 (x\cdot \ln x) \textnormal{d}x = \left[{1\over 2}x^2 \cdot \left(\ln x -{1\over2}\right)\right]_u^1=$$
		$$={1\over 2} \cdot \left(0 -{1\over2}\right)-{1\over 2}u^2 \cdot \left(\ln u -{1\over2}\right)$$
		$$=-{1\over 4}-{1\over 2}u^2 \cdot \left(\ln u -{1\over2}\right)$$
		$$\lim\limits_{u\to0}-{1\over 4}-{1\over 2}u^2 \cdot \left(\ln u -{1\over2}\right)=-{1\over 4}$$
		$$A={1\over 4}~\textnormal{FE}$$
\end{minipage}
\hfill
\begin{minipage}{0.5\textwidth}
	\noindent{\bf Aufgabe 1. h)}\\
	\includegraphics[width=5 cm]{lefz086}		
\end{minipage}

\begin{center}
	\includegraphics[width=7 cm]{../../viecher/endcomic.pdf}
	
	Hier geht es zurück zum \href{https://www.okuyakl.de/math/m12mlnL086/aa086.pdf}{Aufgabenblatt}
\end{center}

\end{document}

